% =============================================================================
% CHAPTER 12: WILLINGNESS TO ACT — Standalone Version
% =============================================================================
% EBF Framework | Version 1.0 | January 2026
% Kompilieren mit: pdflatex chapter_12_standalone.tex
% =============================================================================

\documentclass[11pt,a4paper]{book}

% -----------------------------------------------------------------------------
% PAKETE
% -----------------------------------------------------------------------------
\usepackage[utf8]{inputenc}
\usepackage[T1]{fontenc}
\usepackage[ngerman,english]{babel}
\usepackage{lmodern}

% Mathematik
\usepackage{amsmath,amssymb,amsthm}

% Tabellen
\usepackage{booktabs}
\usepackage{array}
\usepackage{longtable}
\usepackage{colortbl}

% Farben und Boxen
\usepackage[table]{xcolor}
\usepackage{tcolorbox}
\tcbuselibrary{skins,breakable}

% Layout
\usepackage[margin=2.5cm]{geometry}
\usepackage{enumitem}
\usepackage{fancyhdr}
\usepackage{titlesec}

% Hyperlinks
\usepackage{hyperref}
\hypersetup{
    colorlinks=true,
    linkcolor=blue!60!black,
    urlcolor=blue!60!black,
    citecolor=green!50!black
}

% -----------------------------------------------------------------------------
% THEOREME UND DEFINITIONEN
% -----------------------------------------------------------------------------
\theoremstyle{definition}
\newtheorem{definition}{Definition}[section]
\newtheorem{example}{Example}[section]

\theoremstyle{remark}
\newtheorem*{intuition}{Intuition}

% -----------------------------------------------------------------------------
% HEADER/FOOTER
% -----------------------------------------------------------------------------
\pagestyle{fancy}
\fancyhf{}
\fancyhead[LE,RO]{\thepage}
\fancyhead[LO]{Chapter 12: Willingness to Act}
\fancyhead[RE]{EBF Framework}
\renewcommand{\headrulewidth}{0.4pt}

% -----------------------------------------------------------------------------
% TITEL
% -----------------------------------------------------------------------------
\title{
    \vspace{-2cm}
    {\Large EBF — Behavioral Complementarity Model}\\[0.5cm]
    \rule{\textwidth}{1pt}\\[0.5cm]
    {\Huge\bfseries Chapter 12}\\[0.3cm]
    {\LARGE Willingness to Act}\\[0.2cm]
    {\large From Thinking to Deciding}\\[0.5cm]
    \rule{\textwidth}{1pt}
}
\author{
    FehrAdvice \& Partners AG\\[0.2cm]
    \small CORE-READY (AV) — The Seventh Fundamental Question
}
\date{January 2026 — Version 1.0}

% =============================================================================
% DOKUMENT
% =============================================================================
\begin{document}

\maketitle
\thispagestyle{empty}

% -----------------------------------------------------------------------------
% INHALTSVERZEICHNIS
% -----------------------------------------------------------------------------
\tableofcontents
\newpage

% -----------------------------------------------------------------------------
% KAPITEL
% -----------------------------------------------------------------------------
\chapter{Willingness to Act: From Thinking to Deciding}
\label{ch:willingness}

%------------------------------------------------------------------------------
% NAVIGATION BOX
%------------------------------------------------------------------------------
\begin{tcolorbox}[
  title={\textbf{Die 9 Fragen des EBF — Kapitelnavigation}},
  colback=blue!3,
  colframe=blue!50!black,
  fonttitle=\bfseries\sffamily,
  coltitle=white,
  boxrule=0.8pt,
  arc=2mm
]
\small\sffamily
\begin{tabular}{@{}clr@{}}
$\checkmark$ & 0. Kontext und Zielverhalten definieren & [Kapitel 6] \\
$\checkmark$ & 2. WELCHEN NUTZEN bringt das Verhalten? & [Kapitel 10] \\
$\checkmark$ & 3. WIE NIMMT er/sie den Nutzen WAHR? & [Kapitel 11] \\
$\rightarrow$ & \textbf{1. WIE denkt er/sie?} & \textbf{[DIESES KAPITEL]} \\
$\rightarrow$ & \textbf{4. WIE GROSS ist WAX und WELCHE SCHWELLEN?} & \textbf{[DIESES KAPITEL]} \\
$\circ$ & 5. WO in der Behavioral Change Journey? & [Kapitel 13] \\
$\circ$ & 6. WELCHES Behavioral Change Segment? & [Kapitel 14] \\
$\circ$ & 7. WILLINGNESS TO EFFECTIVELY CHANGE? & [Kapitel 15] \\
$\circ$ & 8. EFFEKTIVE WAHRSCHEINLICHKEIT? & [Kapitel 16] \\
\end{tabular}
\end{tcolorbox}

\vspace{1cm}

%------------------------------------------------------------------------------
% CORE-READY BOX
%------------------------------------------------------------------------------
\begin{tcolorbox}[colback=purple!5!white, colframe=purple!75!black,
    title=\textbf{CORE-READY: The Seventh Fundamental Question (7C)}]
\textbf{This chapter answers:} \textit{How ready am I to act?} (Wie handlungsbereit bin ich?)

\textbf{Formal Specification:} Appendix AV (CORE-READY) — 99 axioms, 94+ references

\textbf{The Core Equation:}
\[
\boxed{\text{Action} \Leftrightarrow WAX \geq \theta}
\]

\textbf{Builds on CORE-AWARE (AU):}
\begin{itemize}[nosep]
    \item AWX transforms potential utility into effective utility: $U^* = A(\cdot) \times U_{pot}$
    \item This chapter addresses: Given $U^*$, \textit{will I act}?
    \item The ``gap between knowing and doing'' is the WAX/$\theta$ mechanism
\end{itemize}

\textbf{The Intention-Behavior Gap:} Why knowing $\neq$ doing is explained by threshold dynamics.
\end{tcolorbox}

%%%%%%%%%%%%%%%%%%%%%%%%%%%%%%%%%%%%%%%%%%%%%%%%%%%%%%%%%%%%%%%%%%%%%%%%%%%%%%%
\section{Introduction: The Action Condition}
\label{sec:12-intro}
%%%%%%%%%%%%%%%%%%%%%%%%%%%%%%%%%%%%%%%%%%%%%%%%%%%%%%%%%%%%%%%%%%%%%%%%%%%%%%%

Chapter 11 established \textit{Awareness} (AWX) — what people perceive. But perception alone does not guarantee action. This chapter addresses the next critical question: \textbf{Willingness to Act} (WAX).

\begin{definition}[Willingness to Act]
\textbf{Willingness to Act (WAX)} is the degree to which an individual is prepared to execute a behavior, given their awareness of its consequences. WAX integrates:
\begin{itemize}
\item How the person \textit{thinks} about the decision (thinking style $\varphi$)
\item The perceived utilities across all dimensions ($U^* = AWX \cdot U$)
\item The barriers and facilitators in the decision context
\end{itemize}
\end{definition}

\begin{intuition}
\textbf{Awareness vs. Willingness:}
\begin{itemize}
\item \textbf{Awareness} answers: ``Do I \textit{see} the benefits?''
\item \textbf{Willingness} answers: ``Am I \textit{ready} to act on them?''
\end{itemize}
A person can be fully aware of the benefits of exercise but still unwilling to go to the gym. The gap between knowing and doing is where Willingness operates.
\end{intuition}

%%%%%%%%%%%%%%%%%%%%%%%%%%%%%%%%%%%%%%%%%%%%%%%%%%%%%%%%%%%%%%%%%%%%%%%%%%%%%%%
\section{Question 1: How Does the Person Think? ($\varphi$)}
\label{sec:12-phi}
%%%%%%%%%%%%%%%%%%%%%%%%%%%%%%%%%%%%%%%%%%%%%%%%%%%%%%%%%%%%%%%%%%%%%%%%%%%%%%%

The first question this chapter answers is fundamental: \textbf{How does the person think about multi-dimensional decisions?}

\subsection{The Thinking Style Parameter $\varphi$}

\begin{equation}
\boxed{\varphi \in [0, 1]}
\end{equation}

\begin{intuition}
\textbf{What is $\varphi$?}

When facing a decision with multiple dimensions (financial, social, environmental, identity...), people think in fundamentally different ways:

\begin{center}
\begin{tabular}{cp{10cm}}
\toprule
$\varphi \to 1$ & \textbf{``Connected Thinking'' (Complementary)}\\
& ``I only act if EVERYTHING is right''\\
& The weakest dimension dominates the decision\\
& Like a chain: only as strong as the weakest link\\
& Example: ``I won't buy unless price AND environment AND social proof all check out''\\[8pt]

$\varphi \to 0$ & \textbf{``Separated Thinking'' (Additive)}\\
& ``I look at the total benefit''\\
& All dimensions add up; strengths compensate for weaknesses\\
& Example: ``Great financial savings can make up for weak environmental benefits''\\
\bottomrule
\end{tabular}
\end{center}
\end{intuition}

\subsection{$\varphi$-Decomposition}

The thinking style is not fixed — it depends on both the person and the context:

\begin{equation}
\boxed{\varphi = \varphi_0 + \Delta\varphi(\Psi)}
\end{equation}

where:
\begin{itemize}
\item $\varphi_0$ = baseline thinking style (personality trait, relatively stable)
\item $\Delta\varphi(\Psi) = \sum_j \gamma_j \cdot \Psi_j$ = context-induced shift
\end{itemize}

\begin{intuition}
\textbf{What shifts $\varphi$?}

\begin{center}
\begin{tabular}{lcp{6cm}}
\toprule
\textbf{Context} & \textbf{Effect on $\varphi$} & \textbf{Why}\\
\midrule
High uncertainty ($\Psi_{UNS}$) & $\downarrow$ & Fall back to self-interest\\
Strong group identity ($\Psi_{SEG}$) & $\uparrow$ & Norms become binding\\
Salient identity ($\Psi_{IDN}$) & $\uparrow$ & ``Who am I?'' dominates\\
Social observation ($\Psi_{SOC}$) & $\uparrow$ & Others are watching\\
\bottomrule
\end{tabular}
\end{center}

\textbf{Example:} Anna normally thinks additively ($\varphi_0 = 0.35$). But when her neighbor mentions ``everyone on our street is getting heat pumps,'' her $\varphi$ shifts upward (+0.20) because social norms become salient.
\end{intuition}

%%%%%%%%%%%%%%%%%%%%%%%%%%%%%%%%%%%%%%%%%%%%%%%%%%%%%%%%%%%%%%%%%%%%%%%%%%%%%%%
\section{Question 4a: How Large is the Willingness? (WAX)}
\label{sec:12-wax}
%%%%%%%%%%%%%%%%%%%%%%%%%%%%%%%%%%%%%%%%%%%%%%%%%%%%%%%%%%%%%%%%%%%%%%%%%%%%%%%

\subsection{The WAX Formula}

\begin{equation}
\boxed{WAX = \varphi \cdot \min_i\{U_i^*\} + (1-\varphi) \cdot \sum_i U_i^*}
\end{equation}

where $U_i^* = AWX_i \cdot U_i$ is the \textit{effective utility} — the utility the person actually perceives.

\begin{intuition}
\textbf{What does this formula mean?}

The formula is a weighted combination of two ways of evaluating options:

\begin{enumerate}
\item \textbf{Complementary logic} ($\min$): The weakest dimension determines value\\
$\rightarrow$ ``A meal is only as good as its worst dish''

\item \textbf{Additive logic} ($\sum$): All dimensions contribute to total value\\
$\rightarrow$ ``A great dessert can compensate for mediocre soup''
\end{enumerate}

$\varphi$ determines the weight between these two logics.
\end{intuition}

\subsection{Worked Example: Anna's WAX Calculation}

\begin{example}[Anna's Willingness]
Anna is considering a heat pump. Her values:

\begin{center}
\begin{tabular}{lcccc}
\toprule
\textbf{Dimension} & $U$ (objective) & $AWX$ & $U^* = AWX \cdot U$ \\
\midrule
Individual (INU) & 0.7 & 0.5 & 0.35 \\
Collective (KNU) & 0.6 & 0.5 & 0.30 \\
Identity (IDU) & 0.5 & 0.5 & 0.25 \\
\bottomrule
\end{tabular}
\end{center}

With $\varphi = 0.35$:
\begin{align*}
\min\{U^*\} &= \min\{0.35, 0.30, 0.25\} = 0.25\\
\sum U^* &= 0.35 + 0.30 + 0.25 = 0.90\\
WAX &= 0.35 \cdot 0.25 + 0.65 \cdot 0.90\\
&= 0.088 + 0.585 = \mathbf{0.67}
\end{align*}
\end{example}

%%%%%%%%%%%%%%%%%%%%%%%%%%%%%%%%%%%%%%%%%%%%%%%%%%%%%%%%%%%%%%%%%%%%%%%%%%%%%%%
\section{Question 4b: What Are the Thresholds? ($\theta$)}
\label{sec:12-theta}
%%%%%%%%%%%%%%%%%%%%%%%%%%%%%%%%%%%%%%%%%%%%%%%%%%%%%%%%%%%%%%%%%%%%%%%%%%%%%%%

Willingness alone is not sufficient. It must \textit{exceed a threshold} for action to occur.

\subsection{The Threshold Formula}

\begin{equation}
\boxed{\theta = \theta_{base}(d) + \sum_j \delta_j \cdot \Psi_j}
\end{equation}

\begin{intuition}
\textbf{What determines the threshold?}

\begin{enumerate}
\item $\theta_{base}$: \textbf{Base difficulty} depending on decision type

\begin{center}
\begin{tabular}{lcp{5cm}}
\toprule
\textbf{Decision Type} & $\theta_{base}$ & \textbf{Example}\\
\midrule
Continue habit & 0.2 & Keep using gas heating\\
Small adjustment & 0.5 & Switch electricity provider\\
Moderate change & 0.8 & Buy an e-bike\\
Major investment & 1.2 & Install heat pump\\
Life-changing & 1.5+ & Move abroad for climate\\
\bottomrule
\end{tabular}
\end{center}

\item $\sum \delta_j \cdot \Psi_j$: \textbf{Context adjustments}
\begin{itemize}
\item High uncertainty $\rightarrow$ threshold goes UP (harder to decide)
\item Good subsidies ($\Psi_{REG}$) $\rightarrow$ threshold goes DOWN (easier)
\item Social proof ($\Psi_{SOC}$) $\rightarrow$ threshold goes DOWN
\end{itemize}
\end{enumerate}
\end{intuition}

%%%%%%%%%%%%%%%%%%%%%%%%%%%%%%%%%%%%%%%%%%%%%%%%%%%%%%%%%%%%%%%%%%%%%%%%%%%%%%%
\section{The Critical Comparison: WAX vs. $\theta$}
\label{sec:12-comparison}
%%%%%%%%%%%%%%%%%%%%%%%%%%%%%%%%%%%%%%%%%%%%%%%%%%%%%%%%%%%%%%%%%%%%%%%%%%%%%%%

The entire EBF framework culminates in one comparison:

\begin{center}
\fbox{\parbox{0.6\textwidth}{
\centering
\Large
$WAX \gtrless \theta$
\\[6pt]
\normalsize
\textit{Does willingness exceed the threshold?}
}}
\end{center}

\subsection{The Three Regimes}

\begin{table}[htbp]
\centering
\caption{The Three Decision Regimes}
\begin{tabular}{cccp{5cm}}
\toprule
\textbf{Regime} & \textbf{Condition} & \textbf{P(Action)} & \textbf{Plain Language}\\
\midrule
\cellcolor{red!15} Below & $WAX < \theta$ & $< 50\%$ & ``I don't want it enough''\\
\cellcolor{yellow!15} Critical & $WAX \approx \theta$ & $\approx 50\%$ & ``I'm on the fence''\\
\cellcolor{green!15} Above & $WAX > \theta$ & $> 50\%$ & ``Yes, I'll do it''\\
\bottomrule
\end{tabular}
\end{table}

\subsection{The Gap: $\Delta = WAX - \theta$}

\begin{equation}
\boxed{\Delta = WAX - \theta}
\end{equation}

\begin{intuition}
The gap tells you \textbf{how far from action} someone is:

\begin{center}
\begin{tabular}{rcp{6cm}}
$\Delta > 0$ & $\checkmark$ & Ready to act\\
$\Delta < 0$ & $\times$ & Not ready\\
$\Delta \approx 0$ & $\sim$ & Tipping point (Moment of Truth)\\
\end{tabular}
\end{center}

\textbf{Intervention implication:}
\begin{itemize}
\item Small gap ($|\Delta| < 0.2$): Light nudge sufficient
\item Medium gap ($|\Delta| = 0.2$--$0.5$): Substantive intervention needed
\item Large gap ($|\Delta| > 0.5$): Structural change required
\end{itemize}
\end{intuition}

\subsection{Two Paths to Close the Gap}

To move someone from ``no'' to ``yes,'' there are two complementary paths:

\begin{center}
\begin{tabular}{lp{5cm}p{5cm}}
\toprule
& \textbf{Path 1: Increase WAX} & \textbf{Path 2: Decrease $\theta$}\\
& ``Make it more attractive'' & ``Make it easier''\\
\midrule
Levers & $\uparrow \varphi$ (shift mindset) & $\downarrow \theta_{base}$ (simplify decision)\\
& $\uparrow AWX$ (increase awareness) & $\downarrow$ friction (remove barriers)\\
& $\uparrow U$ (enhance benefits) & Change defaults (opt-out vs. opt-in)\\
\midrule
Example & ``Your neighbors got one too!'' & ``One-stop-shop installation''\\
& ``You'll save CHF 2,000/year'' & ``Pre-filled application form''\\
\bottomrule
\end{tabular}
\end{center}

\textbf{Best practice:} Combine both paths for maximum effect.

%%%%%%%%%%%%%%%%%%%%%%%%%%%%%%%%%%%%%%%%%%%%%%%%%%%%%%%%%%%%%%%%%%%%%%%%%%%%%%%
\section{Complete Example: Anna's Journey}
\label{sec:12-anna}
%%%%%%%%%%%%%%%%%%%%%%%%%%%%%%%%%%%%%%%%%%%%%%%%%%%%%%%%%%%%%%%%%%%%%%%%%%%%%%%

\begin{example}[Anna: From 34\% to 54\%]

\textbf{BASELINE (before intervention):}
\begin{align*}
\varphi &= 0.35 \quad \text{(mostly additive)}\\
AWX &= 0.5 \quad \text{(sees only half the benefits)}\\
WAX &= 0.67\\
\theta &= 1.32 \quad \text{(major investment + uncertainty)}\\
\Delta &= 0.67 - 1.32 = \mathbf{-0.65}\\
P &= \sigma(-0.65) \approx \mathbf{34\%}
\end{align*}

\textbf{INTERVENTION:}
\begin{enumerate}
\item Show concrete savings calculation $\rightarrow$ $AWX$: $0.5 \to 0.8$
\item ``The Müllers got one!'' (social norm) $\rightarrow$ $\varphi$: $0.35 \to 0.55$
\item Highlight subsidies $\rightarrow$ $\theta$: $-0.30$
\item Simplified application $\rightarrow$ $\theta$: $-0.26$
\end{enumerate}

\textbf{AFTER INTERVENTION:}
\begin{align*}
\varphi &= 0.55\\
AWX &= 0.8\\
WAX &= 0.90\\
\theta &= 0.76\\
\Delta &= 0.90 - 0.76 = \mathbf{+0.14}\\
P &= \sigma(+0.14) \approx \mathbf{54\%}
\end{align*}

\textbf{Result:} Anna crossed the tipping point. Probability increased by +20 percentage points.
\end{example}

%%%%%%%%%%%%%%%%%%%%%%%%%%%%%%%%%%%%%%%%%%%%%%%%%%%%%%%%%%%%%%%%%%%%%%%%%%%%%%%
\section{Willingness Axioms: Overview}
\label{sec:12-axioms}
%%%%%%%%%%%%%%%%%%%%%%%%%%%%%%%%%%%%%%%%%%%%%%%%%%%%%%%%%%%%%%%%%%%%%%%%%%%%%%%

Appendix AV provides the complete formal specification of all 99+ Willingness axioms. Here we summarize the key categories:

\begin{table}[htbp]
\centering
\caption{Willingness Axiom Categories}
\begin{tabular}{lcp{6cm}}
\toprule
\textbf{Category} & \textbf{Axioms} & \textbf{Key Content}\\
\midrule
Fundamental & WAX-A1 to A18 & Core definitions, $\varphi$, threshold\\
Barrier & WAX-A19 to A40 & What reduces willingness\\
Facilitator & WAX-A41 to A60 & What increases willingness\\
Dynamic & WAX-A61 to A80 & How willingness changes over time\\
Context & WAX-A81 to A99 & $\Psi$-modulation effects\\
Mapping & WAX-A100 to A103 & Market integration, kinship, thresholds\\
\bottomrule
\end{tabular}
\end{table}

%%%%%%%%%%%%%%%%%%%%%%%%%%%%%%%%%%%%%%%%%%%%%%%%%%%%%%%%%%%%%%%%%%%%%%%%%%%%%%%
\section{Chapter Summary}
\label{sec:12-summary}
%%%%%%%%%%%%%%%%%%%%%%%%%%%%%%%%%%%%%%%%%%%%%%%%%%%%%%%%%%%%%%%%%%%%%%%%%%%%%%%

This chapter addressed two of the nine EBF questions:

\begin{tcolorbox}[
  title={\textbf{Chapter 12: Key Takeaways}},
  colback=green!3,
  colframe=green!40!black,
  fonttitle=\bfseries\sffamily
]

\textbf{Question 1: How does the person think?}
\begin{equation*}
\varphi = \varphi_0 + \sum_j \gamma_j \cdot \Psi_j
\end{equation*}
$\varphi \to 1$: ``Everything must be right'' (complementary)\\
$\varphi \to 0$: ``Total benefit counts'' (additive)

\vspace{6pt}
\textbf{Question 4: How large is WAX and what are the thresholds?}
\begin{align*}
WAX &= \varphi \cdot \min\{U^*\} + (1-\varphi) \cdot \sum U^*\\
\theta &= \theta_{base}(d) + \sum_j \delta_j \cdot \Psi_j
\end{align*}

\textbf{The Critical Comparison:}
\begin{equation*}
\Delta = WAX - \theta \quad \Rightarrow \quad P = \sigma(\Delta)
\end{equation*}

\textbf{Two paths to behavioral change:}
\begin{enumerate}
\item Increase WAX (make it attractive)
\item Decrease $\theta$ (make it easy)
\end{enumerate}

\end{tcolorbox}

\vspace{1cm}

\begin{center}
\rule{0.5\textwidth}{0.5pt}\\[0.5cm]
{\Large\bfseries The Core Equation of EBF}\\[0.5cm]
\fbox{\parbox{0.7\textwidth}{
\centering
\LARGE
$\text{Action} \Leftrightarrow WAX \geq \theta$
}}\\[0.5cm]
\rule{0.5\textwidth}{0.5pt}
\end{center}

\vspace{1cm}

\textbf{Next:} Chapter 13 addresses Question 5 — Where in the Behavioral Change Journey?

\textbf{Formal specification:} Appendix AV provides the complete mathematical formalization with 99+ axioms, worked examples, and scientific references.

% =============================================================================
% ENDE
% =============================================================================

\end{document}
